%% =========================================================
\section{Determinant and Choi--Jamio\l{}kowski operator}
\label{sec:channel_determinant_choi_bound}
%% =========================================================

%% Source: \cite[Section~6.1, Eq.~(6.27)]{Wolf2012Quantum};
%% Notes/WolfNoteTexSource/ch06_spectral_properties.tex, lines 520--545.

\begin{definition}[Purity of the Choi--Jamio\l{}kowski operator]
    \label{def:choi_purity}
    \lean{choiPurity, choiPurity_nonneg, choiPurity_eq_sum}\leanok
    \uses{def:choi_matrix}
    For a linear map $T : \MN{D} \to \MN{D}$ with Choi--Jamio\l{}kowski
    operator $\tau$, the \emph{purity} of $\tau$ is $\tr[\tau^\dagger\tau]$,
    the sum
    \begin{align}
        \tr[\tau^\dagger\tau]
        &= \sum_{i_1,j_1,i_2,j_2}
            \bigl|\tau_{(i_1,j_1),(i_2,j_2)}\bigr|^2
        \label{eq:representations_purity_entries}
    \end{align}
    of the squared moduli of the entries of $\tau$; in particular it is
    non-negative.
\end{definition}

\begin{lemma}[Reshuffling]\label{lem:channel_matrix_eq_reshuffled_choi}
    \lean{channelMatrix_apply_eq_choiMatrix, channelMatrix_apply,
        choiMatrix_apply_single}\leanok
    \uses{def:channel_det, def:choi_matrix}
    Let $T : \MN{D} \to \MN{D}$ be a linear map, $\widehat{T}$ its matrix with
    respect to the matrix-unit basis and $\tau$ its Choi--Jamio\l{}kowski
    operator. Then
    \begin{align}
        \widehat{T}_{(i_1,i_2),(j_1,j_2)}
        &= D\,\tau_{(i_1,j_1),(i_2,j_2)}.
        \label{eq:representations_reshuffling}
    \end{align}
    Equivalently $\widehat{T}=D\,\tau^{\Gamma}$, where the involution
    $\tau\mapsto\tau^{\Gamma}$ is defined by
    $\bra{m,n}\tau^{\Gamma}\ket{k,\ell}=\bra{m,k}\tau\ket{n,\ell}$.
    This is \cite[Eq.~(2.22)]{Wolf2012Quantum}.
\end{lemma}

\begin{proof}\leanok
    Both sides are computed on the matrix units $E_{ij}$:
    \begin{align}
        \widehat{T}_{(i_1,i_2),(j_1,j_2)}
        &= \bigl(T(E_{j_1j_2})\bigr)_{i_1i_2},
        \label{eq:representations_channel_matrix_entries}\\
        \tau_{(i_1,i_2),(j_1,j_2)}
        &= \frac1D\bigl(T(E_{i_2j_2})\bigr)_{i_1j_1}.
        \label{eq:representations_choi_entries}
    \end{align}
    Substituting $(i_1,j_1)$ for the row pair and $(i_2,j_2)$ for the column
    pair in (\ref{eq:representations_choi_entries}) gives
    $\frac1D\bigl(T(E_{j_1j_2})\bigr)_{i_1i_2}$, which by
    (\ref{eq:representations_channel_matrix_entries}) is
    $\frac1D\widehat{T}_{(i_1,i_2),(j_1,j_2)}$.
\end{proof}

\begin{lemma}[Purity and the matrix representation]
    \label{lem:choi_purity_eq_hilbert_schmidt}
    \lean{trace_conjTranspose_mul_self_channelMatrix}\leanok
    \uses{def:channel_det, def:choi_purity}
    With the notation of Lemma~\ref{lem:channel_matrix_eq_reshuffled_choi},
    \begin{align}
        \tr[\widehat{T}^\dagger\widehat{T}]
        &= D^2\,\tr[\tau^\dagger\tau].
        \label{eq:representations_purity_factor}
    \end{align}
    For $D\geq1$ this is \cite[Eq.~(6.28)]{Wolf2012Quantum}, which states it in
    the divided form $\tr[\tau^\dagger\tau]=D^{-2}\tr[\widehat{T}^\dagger
    \widehat{T}]$.
\end{lemma}

\begin{proof}\leanok
    \uses{lem:channel_matrix_eq_reshuffled_choi}
    The trace $\tr[A^\dagger A]$ is the sum of the squared moduli of the
    entries of $A$, so (\ref{eq:representations_reshuffling}) gives
    \begin{align}
        \tr[\widehat{T}^\dagger\widehat{T}]
        &= \sum_{i_1,i_2,j_1,j_2}
            \bigl|\widehat{T}_{(i_1,i_2),(j_1,j_2)}\bigr|^2
         = D^2\sum_{i_1,i_2,j_1,j_2}
            \bigl|\tau_{(i_1,j_1),(i_2,j_2)}\bigr|^2.
        \label{eq:representations_purity_entry_sum}
    \end{align}
    Reshuffling permutes the index quadruples $(i_1,i_2,j_1,j_2)$ without
    repetition, so the remaining sum is $\tr[\tau^\dagger\tau]$ by
    (\ref{eq:representations_purity_entries}).
\end{proof}

\begin{lemma}[Determinant bound by the Hilbert--Schmidt norm]
    \label{lem:det_sq_le_hilbert_schmidt_pow}
    \lean{Matrix.pow_card_mul_norm_det_sq_le_trace_conjTranspose_mul_self_pow}
    \leanok
    \uses{lem:trace_det_amgm}
    For a square complex matrix $A$ whose index set has $n$ elements,
    \begin{align}
        n^{n}\,|\det A|^2 &\leq \tr[A^\dagger A]^{n}.
        \label{eq:representations_det_sq_hs_bound}
    \end{align}
\end{lemma}

\begin{proof}\leanok
    \uses{lem:trace_det_amgm}
    The matrix $A^\dagger A$ is positive semidefinite and
    $\det(A^\dagger A)=|\det A|^2$, so
    (\ref{eq:representations_trace_det_amgm}) applied to $M=A^\dagger A$ is the
    asserted bound. In the singular values $s_1,\dots,s_n$ of $A$ it reads
    $n^n\prod_i s_i^2\leq\bigl(\sum_i s_i^2\bigr)^n$.
\end{proof}

\begin{theorem}[Determinant and Choi--Jamio\l{}kowski operator]
    \label{thm:channelDet_le_choi_purity}
    \lean{norm_channelDet_le_choiPurity_rpow,
        norm_channelDet_sq_le_choiPurity_pow}\leanok
    \uses{def:channel_det, def:choi_purity}
    Let $T:\MN{D}\to\MN{D}$ be a linear map and $\tau$ the corresponding
    Choi--Jamio\l{}kowski operator. Then
    \begin{align}
        |\det T| &\le \tr[\tau^\dagger\tau]^{D^2/2}.
        \label{eq:representations_det_choi_bound}
    \end{align}
    No positivity, trace preservation or unitality is assumed. This is
    \cite[Eq.~(6.27)]{Wolf2012Quantum}.
\end{theorem}

\begin{proof}\leanok
    \uses{lem:choi_purity_eq_hilbert_schmidt, lem:det_sq_le_hilbert_schmidt_pow}
    The matrix $\widehat{T}$ is indexed by the $D^2$ matrix units and
    $\det T=\det\widehat{T}$, so (\ref{eq:representations_det_sq_hs_bound}) with
    $n=D^2$, followed by (\ref{eq:representations_purity_factor}), gives
    \begin{align}
        (D^2)^{D^2}\,|\det T|^2
        &\le \tr[\widehat{T}^\dagger\widehat{T}]^{D^2}
        = (D^2)^{D^2}\,\tr[\tau^\dagger\tau]^{D^2}.
        \notag
    \end{align}
    Cancelling the positive factor $(D^2)^{D^2}$ gives
    $|\det T|^2\le\tr[\tau^\dagger\tau]^{D^2}$, and taking square roots gives
    the stated bound.
\end{proof}
